\section{관련 연구}
\label{sec:related}

\subsection{경기 중 결과와 위험의 예측}
\label{sec:related:esports}

% P1: 경기 결과 예측에서 경기 중 위험 예측으로, 배틀로얄 선행 설계와의 차이
\rev{PUBG의 결과 예측은 경기 종료 후의 요약 통계로 최종 순위를 추정하는 접근에서 출발하였고~\cite{KagglePUBG,LeeLee25}, 다른 e스포츠에서는 경기 중 정보로 승리나 개인의 짧은 구간 사망을 예측하였으며~\cite{Hodge21,Katona19,Marshall22}, \rev{농구에서는 선수 추적 자료로 공격 결과의 기대값을 실시간으로 추정하였다~\cite{Cervone16}.}
배틀로얄에서는 Sagong 등이 단계별 활동 점수로 생존을~\cite{Sagong25a}, 전략 시퀀스로 승리를 예측하였고~\cite{Sagong25b}, KRAFTON의 예측기는 생존분석으로 팀의 실시간 우승확률을 추정하며 그 공개 구현에는 이산시간 hazard 등의 생존모형 기준선도 들어 있다~\cite{KRAFTON}.}
본 연구는 두 설계를 같은 조건에 맞추어 설계 요인의 영향을 구분한다(표~\ref{tab:compare}).

\subsection{생존분석, 생태학적 관점, 다중 에이전트 모델}
\label{sec:related:survival}

% P2: 상대위험과 이산시간 hazard
Cox 비례위험모형은 기저위험을 특정하지 않고 공변량의 상대위험을 추정하며~\cite{Cox72}, 이산시간 생존모형은 구간별 조건부 사건확률을 구간별 기저위험과 함께 모형화한다~\cite{SingerWillett93}.
이 구조는 신경망 생존모형으로 확장되었고~\cite{Gensheimer19}, 예측 시점의 생존자에 대해 이후 위험을 추정하는 landmarking은 동적 예측을 구성하는 방법이다~\cite{vanHouwelingen07}.
상대위험 점수도 기저위험을 결합하면 구간별 확률을 산출할 수 있으므로~\cite{PrenticeGloeckler78}, \rev{본 연구는 기저위험을 결합한 뒤에도 남는 학습 목적의 차이를 따로 측정한다.}

% P3: 생태학적 관점과 모델 구조
생태학은 서식지의 양과 배치가 개체군의 존속과 어떻게 연결되는지 다루어 왔으며~\cite{Fahrig03}, \rev{본 연구는} 이 관점을 위험 공변량의 구성과 해석에만 사용하고, 공변량과 위험의 관계는 경기 자료에서 검증한다.
한편 다중 에이전트 장면에서는 에이전트 축과 시간 축의 attention을 결합하거나~\cite{Ngiam22} 그래프 구조를 attention 편향으로 넣는 방법이 쓰이며\rev{~\cite{Ying21}}, 스포츠에서도 선수 간 관계 그래프가 활용되었다~\cite{Xenopoulos21}.
본 연구는 이러한 구조를 갖춘 모델을 참조 모델로 두고, 입력과 학습 목적을 고정한 채 구조의 조건부 기여를 측정한다.
